% ============================================================
%  REVISION / REVIEWER RESPONSE LETTER
%  Overleaf kullanımı:
%    1. Bu dosyayı ana makale projenize ekleyin.
%    2. xr paketi, ana makalenizin .aux dosyasını okur.
%       "main" yerine kendi ana .tex dosyanızın adını yazın (uzantısız).
%    3. \R{label-adı} ile doğrudan Section/Figure/Table'a atıfta bulunun.
% ============================================================

\documentclass[12pt, a4paper]{article}

% --- Dil & Kodlama ---
\usepackage[T1]{fontenc}
\usepackage[utf8]{inputenc}
\usepackage[english]{babel}

% --- Çapraz referans: ana makalenin .aux dosyasını oku ---
\usepackage{xr}
\externaldocument[main-]{main}   % <-- "main" = ana .tex dosyanızın adı

% --- Yardımcı paketler ---
\usepackage{hyperref}            % tıklanabilir linkler (Overleaf PDF'de çalışır)
\usepackage{xcolor}
\usepackage{geometry}
\usepackage{enumitem}
\usepackage{mdframed}            % hakeme yanıt kutuları için
\usepackage{microtype}

\geometry{margin=2.5cm}

% ---- Renk tanımları ----
\definecolor{reviewercolor}{RGB}{0,70,127}    % hakem yorumu -- koyu mavi
\definecolor{authorcolor}{RGB}{30,120,30}     % yazar yanıtı  -- koyu yeşil
\definecolor{changedcolor}{RGB}{180,0,0}      % değişiklik    -- kırmızı

% ---- Kısayol komutları ----
% Ana makaledeki label'lara kolayca atıf için:
%   \R{sec:intro}    -> Section X
%   \Rf{fig:results} -> Figure X
%   \Rt{tab:data}    -> Table X
\newcommand{\R}[1]{Section~\ref{main-#1}}
\newcommand{\Rf}[1]{Figure~\ref{main-#1}}
\newcommand{\Rt}[1]{Table~\ref{main-#1}}
\newcommand{\Rp}[1]{page~\pageref{main-#1}}  % sayfa numarası

% ---- Hakem & yanıt ortamları ----
\newenvironment{reviewercomment}{%
  \begin{mdframed}[linecolor=reviewercolor, linewidth=1.5pt,
    backgroundcolor=reviewercolor!6, innertopmargin=6pt]
  \color{reviewercolor}\textbf{Reviewer Comment:}\normalcolor\par\smallskip
}{%
  \end{mdframed}\medskip
}

\newenvironment{authorresponse}{%
  \begin{mdframed}[linecolor=authorcolor, linewidth=1pt,
    backgroundcolor=authorcolor!5, innertopmargin=6pt]
  \color{authorcolor}\textbf{Author Response:}\normalcolor\par\smallskip
}{%
  \end{mdframed}\bigskip
}

% Makalede yapılan değişikliği vurgulamak için:
\newcommand{\changed}[1]{\textcolor{changedcolor}{\textit{#1}}}

% ============================================================
\begin{document}

% ---- Başlık ----
\begin{center}
  {\LARGE\bfseries Response to Reviewers}\\[0.8em]
  {\large Manuscript: [Makale Başlığınız]}\\[0.4em]
  {\large Journal: [Dergi Adı]}\\[0.4em]
  \today
\end{center}

\bigskip
\noindent
We thank the editors and reviewers for their thorough and constructive
feedback. Below we address each comment point by point.
All changes in the revised manuscript are highlighted in \changed{red}.

\bigskip
\hrule
\bigskip

% ============================================================
%  REVIEWER 1
% ============================================================
\section*{Reviewer 1}

% ---- Yorum 1.1 ----
\subsection*{Comment 1.1}

\begin{reviewercomment}
  [Hakemin yorumunu buraya yapıştırın.]
\end{reviewercomment}

\begin{authorresponse}
  Thank you for this observation. We have revised \R{sec:methodology}
  to clarify the procedure (see \Rp{sec:methodology} in the revised manuscript).
  Specifically, we now explain that \ldots

  % Doğrudan bir Figure'a atıf:
  As shown in \Rf{fig:pipeline}, the updated workflow \ldots

  % Değişikliği vurgula:
  \changed{``We revised the relevant paragraph in Section~\ref{main-sec:methodology}
  as follows: \ldots''}
\end{authorresponse}

% ---- Yorum 1.2 ----
\subsection*{Comment 1.2}

\begin{reviewercomment}
  [İkinci hakemin yorumunu buraya yapıştırın.]
\end{reviewercomment}

\begin{authorresponse}
  We agree. In the revised version, \Rt{tab:results} has been updated to
  include the additional statistics requested (\Rp{tab:results}).
  \changed{``A new column reporting effect sizes has been added to
  Table~\ref{main-tab:results}.''}
\end{authorresponse}

% ============================================================
%  REVIEWER 2
% ============================================================
\section*{Reviewer 2}

\subsection*{Comment 2.1}

\begin{reviewercomment}
  [Hakemin yorumunu buraya yapıştırın.]
\end{reviewercomment}

\begin{authorresponse}
  This is addressed in \R{sec:discussion}.  We have added a new paragraph
  that discusses the limitations raised by the reviewer.
  \changed{``We added the following paragraph to \R{sec:discussion}: \ldots''}
\end{authorresponse}

% ============================================================
%  EDITOR COMMENTS (varsa)
% ============================================================
\section*{Editor Comments}

\subsection*{Comment E.1}

\begin{reviewercomment}
  [Editörün yorumunu buraya yapıştırın.]
\end{reviewercomment}

\begin{authorresponse}
  As suggested, we have \ldots
\end{authorresponse}

\end{document}
